% foundations.tex — five foundational propositions of math-coding.

\section{Foundations}

Five foundational propositions ground the convention. Each is
instantiated as a packet under \texttt{math/foundations/}.

\subsection{curry-howard: decision is a triple}

\begin{theorem}[curry-howard.triple]
A decision in math-coding is structurally a triple $(\Prop, \Code, \Witness)$,
and the structural correspondence between proposition and code is
verified through witness.
\end{theorem}

\begin{proof}
Curry-Howard isomorphism states that proposition is a type, code is a
term, and inhabitation is a proof. In software engineering, proposition
becomes a statement about a decision, code its realisation, and
witness the git commit binding the proposition to that realisation.
Without the triple, proposition is documentation without consequences:
its change cannot be tracked. With the triple, proposition and code
become a verifiable pair; any code change either updates the
proposition or triggers an explicit supersession.
\end{proof}

\begin{theorem}[Triple in OCaml]
The OCaml record \texttt{decision} mirrors the triple directly:
\begin{verbatim}
type decision = {
  proposition : string;     (* type *)
  code        : string option; (* term *)
  witness     : Sha.t option;  (* inhabitation proof *)
}
\end{verbatim}
\end{theorem}

\subsection{temporal: lifecycle is computed}

\begin{theorem}[temporal.lifecycle]
The lifecycle of a decision $L$ is a function of the decision and git
history; $L$ is computed, not stored.
\end{theorem}

\begin{proof}
Decisions evolve. A proposition true at commit $A$ may be false at commit
$B$. Storing lifecycle as a stored field admits lying: a developer
changes code but forgets to update the \texttt{state} field. Computing
lifecycle from git history makes drift detection automatic and
falsifiable---a lie about state becomes mechanically impossible,
because state is a consequence of history, not an editable field.
\end{proof}

\begin{theorem}[Lifecycle function]
\begin{align*}
L : \Dec \times \mathrm{Repo} &\to \Lifecycle \\
L(d, r) = \;& \textsc{Draft}   &&\text{if } d.\mathtt{witness} = \bot \\
           &\;\mid\; \textsc{Drift} &&\text{if } d.\mathtt{witness} = \text{sha} \;\wedge\; \mathrm{prop}(r, \text{sha}, d.\mathtt{name}) \neq d.\mathtt{proposition} \\
           &\;\mid\; \textsc{Stale} &&\text{if } d.\mathtt{witness} = \text{sha} \;\wedge\; \mathrm{prop}(r, \text{sha}, d.\mathtt{name}) = d.\mathtt{proposition} \;\wedge\; \text{files changed since sha} \\
           &\;\mid\; \textsc{Applied} &&\text{otherwise}
\end{align*}
\end{theorem}

\subsection{constructive: proof is reproducible}

\begin{theorem}[constructive.proof]
A packet with \texttt{substrate: shell} or \texttt{substrate: pbt} requires
reproducible evidence: re-running the command yields the recorded
exit code.
\end{theorem}

\begin{proof}
Classical logic permits ``proven'' without construction. Constructive
logic requires a term---a program that normalises to the sought value.
In software engineering, the term is a re-runnable command: test,
benchmark, type-check. If the result is not reproducible, the claim is
not proven---it is a hypothesis. This is the only epistemic marker
that triggers runtime action; \texttt{register: fact}, \texttt{hypothesis},
\texttt{judgment}, \texttt{unknown} are declarative stances.
\end{proof}

\begin{theorem}[Reproducible definition]
$$\mathrm{Reproducible} \triangleq \mathrm{Re\!-\!run}(\mathrm{command}) = \mathrm{recorded\_exit}$$
\end{theorem}

\subsection{categorical: supersession is SPO}

\begin{theorem}[categorical.supersession-spo]
Supersession $\preceq$ on decisions is a strict partial order:
irreflexive, asymmetric, transitive.
\end{theorem}

\begin{proof}
Decisions evolve. When a proposition changes, the old decision is not
edited---it is superseded by a new one. The relation ``supersedes''
must be a strict partial order to preserve meaning.
Irreflexivity prevents trivial self-supersession. Asymmetry excludes
cycles. Transitivity composes chains into lineages. Without these
properties, the history of decisions becomes inconsistent: a
reviewer cannot tell which decision is current.
\end{proof}

\begin{theorem}[SPO axioms]
\begin{align*}
\forall a.\ &\neg(a \preceq a)                              &&\text{(irreflexive)} \\
\forall a, b.\ &(a \preceq b) \Rightarrow \neg(b \preceq a) &&\text{(asymmetric)} \\
\forall a, b, c.\ &(a \preceq b) \wedge (b \preceq c) \Rightarrow (a \preceq c) &&\text{(transitive)}
\end{align*}
\end{theorem}

\subsection{motivation: register and why}

\begin{theorem}[motivation.register-and-why]
Every packet declares an epistemic register
(\textsc{Fact}/\textsc{Hypothesis}/\textsc{Judgment}/\textsc{Unknown}) and
a numerical confidence. A convention without motivation is ceremony.
\end{theorem}

\begin{proof}
Decisions are taken with different degrees of confidence. A fact
requires reproduction; a hypothesis is open to refutation; a judgment
is an explicit indication of potential harm (reversibility,
mitigation); an unknown is an honest admission of a boundary. Without
explicit register, epistemic honesty degrades into self-deception:
a developer fills ``fact'' without evidence; an agent fills
``fact'' without verification. Without motivation, proposition
becomes a slogan beyond critique. This is axiom A1 (Care) of
math-coding v0.854, restored in v2.0-Y.
\end{proof}

\begin{theorem}[Register-confidence consistency]
The register constrains the confidence value:
\begin{align*}
\mathtt{register} = \textsc{Fact}       &\;\Rightarrow\; \mathtt{confidence} \in [0.95,\, 1] \\
\mathtt{register} = \textsc{Hypothesis} &\;\Rightarrow\; \mathtt{confidence} \in (0.5,\, 0.95) \\
\mathtt{register} = \textsc{Judgment}   &\;\Rightarrow\; \mathtt{confidence} \in \{0,\, 1\} \\
\mathtt{register} = \textsc{Unknown}    &\;\Rightarrow\; \mathtt{confidence} = 0
\end{align*}
The \texttt{judgment} register additionally obligates a non-empty
\texttt{\#\# Why} section in the packet body, declaring reversibility
and mitigation.
\end{theorem}